\documentclass{article}
\usepackage[utf8]{inputenc}
\usepackage[T2A]{fontenc}
\usepackage[russian]{babel}
\usepackage{amsmath, amssymb, amsthm, amsfonts}
\usepackage{geometry}
\geometry{a4paper, margin=2cm}

\theoremstyle{plain}
\newtheorem{theorem}{Теорема}[section]
\newtheorem{lemma}[theorem]{Лемма}
\newtheorem{corollary}[theorem]{Следствие}
\newtheorem{definition}[theorem]{Определение}
\newtheorem{example}[theorem]{Пример}

\title{Конспект ответов на экзаменационные вопросы \\ по методам оптимизации}
\author{}
\date{}

\begin{document}

\maketitle

\section{Математическое программирование. Типы экстремумов функций многих переменных, условия локального экстремума, метод множителей Лагранжа, их интерпретация. Основные понятия выпуклого программирования. Седловые точки. Функция Лагранжа.}

\subsection{Определение математического программирования}
\textbf{Математическое программирование} — раздел математики, изучающий методы решения задач нахождения экстремумов (минимума или максимума) функций на множествах, задаваемых ограничениями (равенствами и неравенствами). Общая задача математического программирования имеет вид:
\[
f(x) \to \min \;(\text{или } \max),\quad x \in X \subset \mathbb{R}^n,
\]
где $f(x)$ — целевая функция, $X$ — допустимое множество, определяемое системой ограничений. В зависимости от свойств функций и ограничений выделяют линейное, нелинейное, выпуклое, квадратичное, дискретное и другие виды программирования.

\subsection{Типы экстремумов}
\begin{definition}
Точка $x^* \in X$ называется точкой \textbf{глобального минимума} функции $f$ на $X$, если $f(x^*) \le f(x)$ для всех $x \in X$.
\end{definition}
\begin{definition}
Точка $x^* \in X$ называется точкой \textbf{локального минимума}, если существует $\varepsilon > 0$ такое, что $f(x^*) \le f(x)$ для всех $x \in X \cap U_\varepsilon(x^*)$, где $U_\varepsilon(x^*)$ — $\varepsilon$-окрестность.
\end{definition}
Аналогично определяются глобальный и локальный максимумы. В зависимости от наличия ограничений различают \textbf{безусловный} экстремум (когда $X = \mathbb{R}^n$) и \textbf{условный} (когда $X$ задаётся ограничениями). Точка экстремума может быть стационарной (если $f$ дифференцируема) или граничной.

\subsection{Необходимые и достаточные условия локального экстремума}
Пусть $f: \mathbb{R}^n \to \mathbb{R}$ непрерывно дифференцируема в окрестности $x^*$.

\begin{theorem}[Необходимое условие (Ферма)]
Если $x^*$ — точка локального экстремума, то $\nabla f(x^*) = 0$.
\end{theorem}
\begin{proof}
Рассмотрим произвольное направление $d \in \mathbb{R}^n$, $\|d\|=1$, и функцию одной переменной $\varphi(t) = f(x^* + t d)$. При $t=0$ функция $\varphi$ имеет локальный экстремум, следовательно, $\varphi'(0) = \langle \nabla f(x^*), d \rangle = 0$. Поскольку $d$ произвольный, $\nabla f(x^*) = 0$.
\end{proof}

Пусть $f$ дважды непрерывно дифференцируема в окрестности $x^*$ и $\nabla f(x^*) = 0$. Рассмотрим матрицу Гессе $H(x^*) = \left(\frac{\partial^2 f}{\partial x_i \partial x_j}(x^*)\right)$.

\begin{theorem}[Достаточное условие]
Если квадратичная форма $\langle H(x^*)y, y \rangle$ строго положительно определена, то $x^*$ — точка локального минимума. Если строго отрицательно определена, то $x^*$ — точка локального максимума. Если форма неопределена, то экстремума нет.
\end{theorem}
\begin{proof}
Из разложения Тейлора:
\[
f(x^*+h) = f(x^*) + \frac{1}{2}\langle H(x^*)h, h\rangle + o(\|h\|^2).
\]
При положительной определённости $H(x^*)$ существует $\lambda > 0$ такое, что $\langle H(x^*)h, h\rangle \ge \lambda \|h\|^2$ для всех $h$. Тогда для достаточно малых $h$ имеем $f(x^*+h) - f(x^*) \ge \frac{\lambda}{4}\|h\|^2 > 0$, что означает локальный минимум. Аналогично для отрицательной определённости. Если форма неопределена, найдутся направления, по которым приращение положительно и отрицательно, значит, экстремума нет.
\end{proof}

Для проверки знакоопределённости применяется критерий Сильвестра: форма положительно определена $\iff$ все главные миноры матрицы Гессе положительны.

\subsection{Метод множителей Лагранжа}
Рассмотрим задачу с ограничениями-равенствами:
\[
f(x) \to \min,\quad \varphi_j(x)=0,\ j=1,\dots,m \quad (m \le n).
\]
Предполагаем, что функции непрерывно дифференцируемы и в точке экстремума $x^*$ матрица Якоби $\left(\frac{\partial \varphi_j}{\partial x_i}\right)$ имеет ранг $m$ (условие регулярности).

\begin{definition}
Функция Лагранжа: $L(x, \lambda) = f(x) + \sum_{j=1}^m \lambda_j \varphi_j(x)$, где $\lambda_j$ — множители Лагранжа.
\end{definition}

\begin{theorem}[Необходимые условия экстремума]
Если $x^*$ — точка экстремума задачи, то существуют такие $\lambda_1^*,\dots,\lambda_m^*$, что
\[
\frac{\partial L}{\partial x_i}(x^*,\lambda^*) = 0,\ i=1,\dots,n,\quad \frac{\partial L}{\partial \lambda_j}(x^*,\lambda^*) = \varphi_j(x^*) = 0.
\]
\end{theorem}
\begin{proof}[Идея доказательства]
При выполнении условия регулярности из теоремы о неявных функциях можно выразить $m$ переменных, например $x_1,\dots,x_m$, через остальные $x_{m+1},\dots,x_n$ в окрестности $x^*$. Тогда задача сводится к безусловной минимизации функции $\tilde{f}(x_{m+1},\dots,x_n) = f(x_1(x_{m+1},\dots,x_n),\dots,x_n)$. Условия стационарности для $\tilde{f}$ дают равенства, которые после применения правила дифференцирования сложной функции приводят к существованию множителей $\lambda_j^*$ и условиям $\frac{\partial L}{\partial x_i}=0$.
\end{proof}

\subsection{Экономическая интерпретация}
В задачах оптимизации ресурсов $\lambda_j$ интерпретируются как \textbf{теневые цены} (объективно обусловленные оценки). Они показывают, на сколько изменится оптимальное значение целевой функции при единичном увеличении ресурса $b_j$ (если ограничение имеет вид $\varphi_j(x) \le b_j$).

\subsection{Основные понятия выпуклого программирования}
\begin{definition}
Множество $X \subset \mathbb{R}^n$ выпукло, если $\forall x,y \in X$, $\forall \alpha \in [0,1]$: $\alpha x + (1-\alpha)y \in X$.
\end{definition}
\begin{definition}
Функция $f: X \to \mathbb{R}$ выпукла, если $\forall x,y \in X$, $\forall \alpha \in (0,1)$: $f(\alpha x + (1-\alpha)y) \le \alpha f(x) + (1-\alpha) f(y)$.
\end{definition}
Свойства: выпуклая функция непрерывна во внутренних точках; для дифференцируемой $f$ выпуклость эквивалентна $f(y) \ge f(x) + \langle \nabla f(x), y-x \rangle$; для дважды дифференцируемой — положительной полуопределённости гессиана. Множество $\{x \in X \mid f(x) \le \lambda\}$ выпукло.

\subsection{Седловые точки и функция Лагранжа}
Рассматривается задача выпуклого программирования:
\[
f(x) \to \min,\quad \varphi_j(x) \le 0,\ j=1,\dots,m,
\]
где $f$ и $\varphi_j$ выпуклы.
Функция Лагранжа: $L(x, \lambda) = f(x) + \sum_{j=1}^m \lambda_j \varphi_j(x)$, $\lambda_j \ge 0$.

\begin{definition}
Пара $(x^*, \lambda^*)$, $\lambda^* \ge 0$, называется \textbf{седловой точкой}, если
\[
L(x^*, \lambda) \le L(x^*, \lambda^*) \le L(x, \lambda^*) \quad \forall x,\ \forall \lambda \ge 0.
\]
\end{definition}
Седловая точка даёт оптимальное решение прямой и двойственной задач: $x^*$ — решение прямой, $\lambda^*$ — решение двойственной.

\section{Теорема Куна – Таккера и ее геометрическая интерпретация. Современные методы градиентной оптимизации.}

\subsection{Теорема Куна – Таккера}
Рассмотрим задачу выпуклого программирования в дифференцируемой форме. Пусть $x^*$ — допустимая точка, $I(x^*) = \{j \mid \varphi_j(x^*) = 0\}$.

\begin{theorem}[Куна – Таккера]
Если выполнено условие регулярности Слейтера (существует $x$ с $\varphi_j(x) < 0$ для всех $j$), то $x^*$ — оптимальное решение тогда и только тогда, когда существуют такие $\lambda_j^* \ge 0$, $j \in I(x^*)$, что
\[
-\nabla f(x^*) = \sum_{j \in I(x^*)} \lambda_j^* \nabla \varphi_j(x^*). \tag{1}
\]
Для задачи с линейными ограничениями условие Слейтера не требуется.
\end{theorem}
\textit{Доказательство (необходимость).} Из теории двойственности следует, что в точке оптимума градиент целевой функции лежит в конусе, порождённом градиентами активных ограничений, и коэффициенты $\lambda_j^*$ неотрицательны. Для выпуклых задач это условие также достаточно.

\textit{Геометрическая интерпретация:} Условие (1) означает, что направление наискорейшего убывания целевой функции ($-\nabla f(x^*)$) является неотрицательной линейной комбинацией градиентов активных ограничений. Если бы существовало направление $d$, для которого $\langle \nabla f(x^*), d \rangle < 0$ и $\langle \nabla \varphi_j(x^*), d \rangle \le 0$ для всех активных $j$, то, двигаясь вдоль $d$, можно было бы уменьшить $f$ и не нарушить ограничений (с точностью до первого порядка). Условие (1) исключает такое направление.

\subsection{Современные методы градиентной оптимизации}
\subsubsection{Градиентный спуск}
Итерационная формула: $x^{k+1} = x^k - h_k \nabla f(x^k)$. Для выпуклых функций с $L$-липшицевым градиентом ($\|\nabla f(y)-\nabla f(x)\|_2 \le L\|y-x\|_2$) выбор $h = 1/L$ даёт
\[
f(\bar{x}^N) - f(x^*) \le \frac{L R^2}{2N},\quad R = \|x^0 - x^*\|_2,\quad \bar{x}^N = \frac{1}{N}\sum_{k=1}^N x^k.
\]
Для $\mu$-сильно выпуклых функций: $f(x^{N+1})-f(x^*) \le (1-\mu/L)^N (f(x^0)-f(x^*))$.

\subsubsection{Ускоренный градиентный метод (Нестерова)}
Двухшаговая схема:
\[
y^{k+1} = x^k + \frac{k}{k+3}(x^k - x^{k-1}),\quad x^{k+1} = y^{k+1} - \frac{1}{L}\nabla f(y^{k+1}).
\]
Оценка: $f(x^N) - f(x^*) \le \frac{2LR^2}{N^2}$, что оптимально для гладких выпуклых задач.

\subsubsection{Стохастический градиентный спуск}
Применяется для $f(x) = \mathbb{E}_{\xi}[f(x,\xi)]$. На итерации используется $\nabla f(x^k, \xi^k)$, где $\xi^k$ — случайная величина. Оценки сходимости: $\mathbb{E}[f(\bar{x}^N)-f(x^*)] = O(1/\sqrt{N})$.

\section{Формулировка задачи линейного программирования (ЛП). Понятия опорного плана и базиса, вырожденность и невырожденность задач ЛП, основные принципы симплекс-метода. Основные теоремы ЛП.}

\subsection{Постановка задачи ЛП}
Каноническая форма:
\[
w(x) = c^T x \to \min,\quad Ax = b,\ x \ge 0,
\]
где $A \in \mathbb{R}^{m \times n}$, $\operatorname{rank} A = m$, $b \ge 0$.

\subsection{Базис, опорный план, вырожденность}
Пусть $B = (A_{\sigma(1)},\dots,A_{\sigma(m)})$ — набор линейно независимых столбцов (базис). Соответствующие переменные $x_B$ — базисные, остальные $x_N$ — свободные. Базисное решение: $x_B = B^{-1}b$, $x_N = 0$. Если $x_B \ge 0$, то оно называется \textbf{опорным планом} (базисным допустимым решением).

\begin{definition}
Опорный план называется \textbf{вырожденным}, если среди его базисных переменных есть нулевые (т.е. $|\{j: x_j > 0\}| < m$). В противном случае план называется \textbf{невырожденным}.
\end{definition}
Вырожденность возникает, когда в системе ограничений есть избыточные связи или когда базисное решение находится на пересечении более чем $m$ гиперплоскостей. Наличие вырожденности может приводить к зацикливанию в симплекс-методе.

\subsection{Симплекс-метод}
Основная идея: переход от одного опорного плана к смежному (отличающемуся одной базисной переменной) с улучшением целевой функции.

Симплекс-таблица строится по базису. Критерий оптимальности: если все оценки $z_{0j} = c_j - c_B B^{-1}A_j \ge 0$, то текущий план оптимален.

Алгоритм:
\begin{enumerate}
\item Найти $s$ с $z_{0s} < 0$ (ведущий столбец). Если таких нет — конец (оптимум).
\item Если все $z_{is} \le 0$, то целевая функция неограничена снизу.
\item Иначе выбрать $r = \arg\min_{i: z_{is}>0} \frac{z_{i0}}{z_{is}}$ (ведущая строка). Элемент $z_{rs}$ — ведущий.
\item Преобразовать таблицу (элементарные преобразования Гаусса), перевести $x_s$ в базис, $x_{\sigma(r)}$ — в свободные.
\item Перейти к шагу 1.
\end{enumerate}

\subsection{Основные теоремы ЛП}
\begin{theorem}[Критерий разрешимости]
Задача ЛП разрешима тогда и только тогда, когда допустимое множество непусто и целевая функция ограничена снизу на нём.
\end{theorem}
\begin{theorem}[Существование опорного плана]
Если допустимое множество непусто, то существует опорный план.
\end{theorem}
\begin{theorem}[Оптимальность на опорном плане]
Если задача разрешима, то существует оптимальный опорный план.
\end{theorem}
\begin{theorem}[Двойственности]
Прямая и двойственная задачи либо обе разрешимы (и тогда их оптимальные значения равны), либо обе неразрешимы (одна из-за несовместности, другая из-за неограниченности). Для оптимальных планов $x^*$, $y^*$ выполняются условия дополняющей нежёсткости: $y_i^*(a_i x^* - b_i) = 0$, $(c_j - y^*A_j)x_j^* = 0$.
\end{theorem}

\section{Потоки в сетях. Теорема Форда – Фалкерсона.}
\subsection{Основные определения}
\begin{definition}
Сеть — ориентированный граф $G=(V,E)$ с источником $s$ и стоком $t$. Каждой дуге $(u,v)$ приписана пропускная способность $c(u,v) \ge 0$. Поток $f$ — функция на дугах, удовлетворяющая:
\begin{enumerate}
\item $0 \le f(u,v) \le c(u,v)$;
\item $\sum_{v} f(u,v) = \sum_{v} f(v,u)$ для всех $u \neq s,t$ (сохранение потока).
\end{enumerate}
Величина потока $|f| = \sum_{v} f(s,v) - \sum_{v} f(v,s)$.
\end{definition}

Разрез $(S,T)$: $s \in S$, $t \in T$, $T = V \setminus S$. Пропускная способность разреза $C(S,T) = \sum_{u\in S, v\in T} c(u,v)$.

\subsection{Теорема Форда – Фалкерсона}
\begin{theorem}[Форда – Фалкерсона]
Максимальная величина потока равна минимальной пропускной способности разреза.
\end{theorem}
\begin{proof}
\textit{1. Оценка сверху:} Для любого потока $f$ и любого разреза $(S,T)$ имеем $|f| = \sum_{u\in S, v\in T} f(u,v) - \sum_{u\in S, v\in T} f(v,u) \le \sum_{u\in S, v\in T} c(u,v) = C(S,T)$. Следовательно, максимальный поток не превышает минимальной пропускной способности разреза.

\textit{2. Построение потока, достигающего этой границы:} Алгоритм Форда – Фалкерсона строит увеличивающие пути в остаточной сети. Остаточная сеть содержит прямые дуги с остаточной пропускной способностью $c(u,v)-f(u,v)$ и обратные дуги с $f(u,v)$. Если увеличивающего пути из $s$ в $t$ в остаточной сети нет, то множество $S$ вершин, достижимых из $s$ по ненасыщенным дугам, не содержит $t$. Для этого разреза $(S,T)$ все прямые дуги из $S$ в $T$ насыщены ($f(u,v)=c(u,v)$), а обратные дуги имеют нулевой поток ($f(v,u)=0$). Следовательно, $|f| = \sum_{u\in S, v\in T} c(u,v) = C(S,T)$. Таким образом, поток достигает пропускной способности разреза, и по предыдущему пункту он максимален.

\textit{3. Конечность при целых пропускных способностях:} При целых $c(u,v)$ алгоритм увеличивает поток на целое положительное число на каждом шаге, поэтому за конечное число шагов будет достигнут максимум.
\end{proof}

\section{Динамическое программирование. Примеры задач, решаемых методом динамического программирования.}

\subsection{Общая постановка и принцип оптимальности}
Рассматривается $N$-шаговый процесс. Состояние на $k$-м шаге $x^{(k)}$, управление $u^{(k)}$. Уравнение состояний: $x^{(k)} = f^{(k)}(x^{(k-1)}, u^{(k)})$. Целевая функция: $J = \sum_{k=1}^N J_k(x^{(k-1)}, u^{(k)})$.

\begin{definition}[Принцип оптимальности Беллмана]
Для любого начального состояния оптимальная стратегия на последующих шагах не зависит от того, как система пришла в это состояние, и определяется только самим состоянием.
\end{definition}

Уравнения Беллмана (для максимизации):
\[
B_N(x^{(N-1)}) = \max_{u^{(N)}} J_N(x^{(N-1)}, u^{(N)}),
\]
\[
B_k(x^{(k-1)}) = \max_{u^{(k)}} \left\{ J_k(x^{(k-1)}, u^{(k)}) + B_{k+1}(x^{(k)}) \right\},\quad k = N-1,\dots,1.
\]
Результатом является $B_1(x^{(0)})$ — максимальное значение целевой функции за все $N$ шагов, и условно-оптимальные управления $u^{*(k)}(x^{(k-1)})$. Затем, зная $x^{(0)}$, последовательно находятся $x^{(1)},\dots,x^{(N)}$ и $u^{(1)},\dots,u^{(N)}$.

\subsection{Примеры задач}

\subsubsection{Задача о распределении средств между предприятиями}
Имеется $N$ предприятий, начальные средства $x^{(0)}$. Прибыль от вложения $u$ в $k$-е предприятие $J_k(u)$. Требуется распределить средства так, чтобы суммарная прибыль была максимальной.

Параметр состояния $x^{(k)}$ — остаток средств после $k$ предприятий. Уравнение состояний: $x^{(k)} = x^{(k-1)} - u^{(k)}$. Уравнения Беллмана решаются таблично, начиная с последнего предприятия. Пример из пособия Ларина: при $x^{(0)}=5$ и заданных $J_k(u)$ оптимальное распределение $(1,2,1,1)$ с суммарной прибылью $24$.

\subsubsection{Задача об оптимальном распределении ресурсов между отраслями на $N$ лет}
Две отрасли, начальные средства $x^{(0)}$. Возврат средств: $q_1(u)$ и $q_2(x-u)$. Прибыль $J_1(u)+J_2(x-u)$. Уравнение состояний: $x^{(k)} = q_1(u^{(k)}) + q_2(x^{(k-1)}-u^{(k)})$. Решение рекуррентных соотношений приводит к оптимальной стратегии инвестиций.

\subsubsection{Задача о замене оборудования}
Оборудование эксплуатируется $N$ лет. Затраты на содержание $r(t)$, ликвидная стоимость $g(t)$. Решение: в начале каждого года принимать решение — сохранить или заменить. Параметр состояния — возраст $t$. Уравнения Беллмана минимизируют суммарные затраты. В примере из пособия при $N=5$ оптимальная стратегия: заменить в начале 3-го года.

\end{document}
